% ==============================================================================
% Section 1: Introduction
% ==============================================================================
\section{Introduction}
\label{sec:intro}

% ---- Block 1: Scale and stakes ------------------------------------------------
Mexico's disappearance crisis is among the most severe humanitarian emergencies in the Western Hemisphere.
As of early 2026, the Registro Nacional de Personas Desaparecidas y No Localizadas (RNPDNO) has recorded over 115,000 persons whose whereabouts remain unknown---a figure widely acknowledged to undercount the true scale \citep{mandolessi_2022, guercke_2021}.
In March 2026, the UN Committee on Enforced Disappearances (CED) issued its first-ever referral to the General Assembly under Article~34 of the Convention, finding that enforced disappearances in Mexico constitute crimes against humanity \citep{CED_MEX_art34_2026}.
The decision explicitly notes that attacks occurred ``at different moments and in different parts of the territory''---a spatial and temporal characterization that, until now, has lacked systematic quantitative support at the national scale.
Forced disappearances have been documented across multiple contexts: as a tool of state repression, a byproduct of organized criminal violence, and a consequence of weak institutional capacity to investigate and resolve cases \citep{solar_2021, guercke_2021, rios_2013}.

% ---- Block 2: Spatial criminology lineage --------------------------------------
Spatial analysis of crime has a well-established lineage in quantitative criminology.
\citet{sherman_1989} demonstrated that predatory crime concentrates at specific places, a finding that launched the hot-spot policing paradigm \citep{braga_hotspots_2014}.
\citet{anselin_1995}'s Local Indicators of Spatial Association (LISA) provided the statistical framework for identifying clusters of high and low values while accounting for spatial dependence.
\citet{weisburd_2015}'s ``law of crime concentration'' formalized the empirical regularity that approximately 50\% of crime occurs in 4--6\% of places, a benchmark that has been validated across cities and crime types \citep{weisburd_2004, curman_2015, steenbeek_weisburd_2016}.
The \emph{Journal of Quantitative Criminology} has dedicated special issues to place-based crime analysis \citep{braga_law_2017, andresen_advances_2021}, and concentration measurement has been refined using Gini coefficients and related indices \citep{prieto_curiel_2018, schnell_braga_piza_2017}.

% ---- Block 3: Mexico-specific spatial work -------------------------------------
Spatial analysis of violence in Mexico has focused primarily on homicides.
\citet{vilalta_muggah_2016} and \citet{vilalta_2021} examined homicide patterns in Mexico City using spatial regression and diffusion models, identifying neighborhood-level clustering tied to socioeconomic deprivation and routine activities.
\citet{cadena_garrocho_2019} provided the first national-scale spatial analysis to include both homicides and disappearances at the municipal level for 2006--2017, documenting a northward shift in violence associated with cartel territorial competition---but treated disappearances as a monolithic category.
\citet{atuesta_2024} examined organized crime-related disappearances in three northern states (Durango, Tamaulipas, Coahuila) using case-level data, documenting spatial patterns linked to drug trafficking routes, but their analysis was restricted to three states and did not examine outcome status decomposition.
\citet{aguilar_velazquez_2025} analyzed disappearances in Mexico City using data-driven methods but at a single-city scale.
None of these studies disaggregated disappearances by outcome status or examined the full national territory at monthly resolution.

% ---- Block 4: The gap ---------------------------------------------------------
All prior spatial work on disappearances in Mexico treats them as a single phenomenon.
This is a consequential simplification.
The RNPDNO distinguishes four outcome statuses that reflect fundamentally different processes: Not Located cases represent active search failures---persons whose whereabouts remain unknown; Located Alive cases represent successful recoveries; and Located Dead cases represent lethal outcomes.
If these processes are driven by different actors, institutional responses, or structural conditions, they should produce different spatial signatures.
A municipality that is a hotspot for living recoveries may not be a hotspot for lethal outcomes; a municipality with many unresolved cases may not be a municipality where bodies are found.
These distinctions have direct implications for resource allocation: forensic identification teams should be deployed where lethal outcomes concentrate, while active search and prevention programs should target areas of concentrated unresolved cases---and these may be different places.
The RNPDNO makes outcome-disaggregated analysis possible, but no study has exploited this structure at the national scale.

% ---- Block 5: This paper -------------------------------------------------------
This paper presents the first outcome-disaggregated spatial analysis of disappearances in Mexico at the national scale with monthly resolution.
We analyze \nmunis{} municipalities over \nmonths{} months (\studyperiod{}) across four outcome statuses, executing \nLISAruns{} LISA runs.
Two methodological contributions distinguish this work.
First, we quantify cross-status spatial non-interchangeability using Jaccard similarity coefficients on pairwise HH cluster sets---demonstrating that outcome statuses produce distinct hotspot geographies.
Second, the monthly temporal resolution reveals dynamics invisible to annual analyses, most notably the rapid expansion of HH clusters in Mexico's Centro region.
We further examine whether the CED's qualitative geographic claims about eight named states find support in the quantitative spatial record.
A companion Data in Brief descriptor documents the dataset schema and processing pipeline \citep{escobar_rnpdno_2026}.

% ---- Block 6: Research questions -----------------------------------------------
This paper addresses four research questions:

\begin{enumerate}
  \item \textbf{RQ1 (Concentration):} How geographically concentrated are
        disappearances, and does concentration differ across outcome statuses?
  \item \textbf{RQ2 (Clustering):} Do the four outcome statuses produce distinct
        spatial cluster geographies at the municipal level?
  \item \textbf{RQ3 (Temporal dynamics):} How has the spatial distribution of
        each status evolved from 2015 to 2025, and do regional hotspot trajectories
        diverge by outcome?
  \item \textbf{RQ4 (Sex disaggregation):} Do male and female disappearances
        exhibit different spatial patterns across outcome statuses?
\end{enumerate}

\noindent We additionally examine the alignment between our empirical findings and the CED's geographic characterization of disappearance patterns across eight named states.
